\section{The actual complex, rational and integral chart values}
\label{sec:chart-values}

\status{CV1--CV9 prove the source's
\texttt{ValuationCompatibleChart} for the true column family, at every
coefficient prime, from its supplied compatible charts. CV10 gives the
same value statements for the smaller terminal row space. The chart
and carrier constructions are treated in FC; the present argument
constructs the rational forms and proves exact integral membership.
Its conclusion has no norm, Petersson, period or dilation clause.}

\paragraph{Notation and the actual data.}
Keep the true column datum \(I:\mathrm{TC}(B,X)\), its
\(\mathrm{Ok}\), \(\mathrm{Pins}\), \(\mathrm{Laws}\) and
\(\mathrm{Later}\) hypotheses, and the resulting record
\(D=I.\mathrm{rec7}(l,Z)\). Fix a word \(u\), a twist \(a\), and
a level \(i\le m\). Write \(E/F\) for the CM extension,
\(L=B.\mathrm{LL}\) for the coefficient number field,
\(\sigma:L\hookrightarrow\mathbb C\) for its prescribed embedding,
and \(j_\tau:E\hookrightarrow L\) for the prescribed finite family
of embeddings. Thus the actual complex embeddings are
\(\iota_\tau=\sigma\circ j_\tau\). Put \(n=N_i\) and let
\(\mathcal C\) be the source's finite class set
\((I.\mathrm{stp}(u,a,i)).\mathrm{Cl}\).
The member-forms record is denoted by \(\mathcal M\), its
\(L\)-space by \(V_{\mathrm{form}}\), and its integral lattice
by \(\mathcal L\).

Each class has a coordinate lattice \(H_c\subset E^n\).
SF gives the literal standard Hermitian form
\(h(x,y)=\overline{x}^{\,t}y\) in these coordinates, and
identifies \(H_c\) with the source's decorated lattice \(L_c\).
Its Hermitian dual is
\[
 H_c^\#=\{y:\overline{x}^{\,t}y\in\mathcal O_E
                           \text{ for every }x\in H_c\}.
\]
The decorated lattice called \texttt{Lsh} is exactly \(H_c^\#\).
The local basis matrices below describe these class lattices in
the fixed global coordinates; they do not replace the common
global frame of SF.

\paragraph{CV1: frame functions retain the determinant twist.}
Let \(\mu_\tau=\mathcal M.\mathrm{lam}(i,\tau)\), and put
\[
 c_i=\frac{i(t+1)}2,\qquad
 \lambda_\tau=\mu_\tau-c_i(1,\ldots,1),\qquad
 d_\tau=\operatorname{detShift}(\lambda_\tau).
                                                               \tag{CV.1}
\]
Here \(t\) is odd, so \(c_i=i((t+1)/2)\) is an integer with
no rounding discrepancy. This identifies the source's two
opposite conventions for \texttt{shW}, hence its two literal
definitions of \(\lambda\). Their polynomial weights
\(\nu_\tau=\operatorname{natW}(\lambda_\tau)
=\operatorname{natW}(\mu_\tau)\) agree.

A frame \(z=(z_\tau)_\tau\) consists of invertible complex
matrices. It need not be orthonormal. For invertible matrices
\(A=(A_\tau)_\tau\), define
\[
 \mathcal E_{\lambda,A}(P)(z)
   =\left(\prod_\tau\det(A_\tau z_\tau)^{d_\tau}\right)
                      P((A_\tau z_\tau)_\tau).
                                                               \tag{CV.2}
\]
This is exactly \texttt{fnOf}; negative determinant powers are
evaluated only on invertible matrices. Evaluation in (CV.2) is
injective. Indeed its determinant factor is nonzero, so a
polynomial in its kernel vanishes on every invertible tuple.
For an arbitrary tuple \(X\), substitute
\(X_\tau+s_\tau I\). Each determinant excludes only finitely
many values of \(s_\tau\). A polynomial vanishing for the
remaining independent values is zero; evaluate at all
\(s_\tau=0\). This also proves injectivity after mapping
coefficients through \(\sigma\).

Use the source's polynomial action
\[
 \rho_\lambda(g)P(X)
   =\left(\prod_\tau\det(g_\tau)^{d_\tau}\right)
                              P((g_\tau^tX_\tau)_\tau).
 \qquad
 \mathcal E_{\lambda,1}(\sigma(\rho_\lambda(g)P))
       =\mathcal E_{\lambda,\sigma(g)^t}(\sigma P).
                                                               \tag{CV.3}
\]
Substitution and determinant multiplication prove this identity
and \(\rho_\lambda(gh)=\rho_\lambda(g)\rho_\lambda(h)\).
Left translating a frame function by an invertible tuple
therefore preserves its polynomial representation space.
The coordinate equivalence \(e_Y\) used below has
\((e_Yz)_\tau=z_\tau\) exactly; it introduces neither a scalar
nor a change of coefficient embedding.

\paragraph{CV2: the rational and integral polynomial spaces are exact.}
For each \(\tau\), let
\(\Delta_{\nu_\tau}(\mathbb Z)\) be the lattice generated by
divided-power words on the highest polynomial, as in CS2.
Let \(\Delta_\lambda^{\mathrm{std}}\subset L[X]\) be the
\(\mathcal O_L\)-span of the products
\(\prod_\tau f_\tau(X_\tau)\), with
\(f_\tau\in\Delta_{\nu_\tau}(\mathbb Z)\).
It is the source's \texttt{delStd}.
Let \(V_\lambda\) be the source's coefficient space
\texttt{VV}. In characteristic zero, IW3 and CS2 give
\[
       L\Delta_\lambda^{\mathrm{std}}=V_\lambda.
\]
This is a field-span equality. The integral Weyl lattice
\(\Delta\) is retained throughout.

At a coefficient prime \(w\), put
\(R_w=(\mathcal O_L)_w\).
The prescribed valuation identifies \texttt{intW} with
\(R_w=\{r:|r|_w\le1\}\). Directly from the generators,
\[
 \begin{split}
 \operatorname{RatSpan}(\mu)
       &=\{\sigma P:P\in V_\lambda\},\\
 \operatorname{IntSpan}_w(\mu)
       &=\{\sigma P:P\in
                      (\Delta_\lambda^{\mathrm{std}})_w\}.
 \end{split}                                                   \tag{CV.4}
\]
For the first equality, expand each Weyl factor in its integral
divided-power generators and distribute the finite tensor
products. Their \(L\)-span is \(V_\lambda\).
For the second, a finite sum with coefficients in \(R_w\)
has a common denominator outside \(w\), so it lies in the
localization of \(\Delta_\lambda^{\mathrm{std}}\).
Conversely, if \(sP\) belongs to that lattice and \(s\notin w\),
expand \(sP\) in the same integral tensors and divide their
coefficients by \(s\). All resulting coefficients remain in
\(R_w\). This proves the literal additive-closure descriptions
of both spans, with no integral index discarded.
The weight shift in (CV.1) leaves the generators unchanged;
its determinant factor remains in (CV.2).

\paragraph{CV3: rational complex functions give actual rational forms.}
For the moment take any source step with class set
\(\mathcal C\), stabilizers \(\Gamma_c\), matrices
\(\gamma_s\in\mathrm{GL}_n(E)\), and coefficient-field
characters \(\eta_s\in L^\times\). Its complex member space
consists of tuples \(\phi_c\in
\mathcal E_{\lambda,1}(V_{\lambda,\mathbb C})\) satisfying
\[
 \phi_c(z)=\sigma(\eta_s)\,
                 \phi_c(\iota(\gamma_s^{-1})z)
                       \quad(s\in\Gamma_c).                  \tag{CV.5}
\]
Its rational forms are tuples \(P_c\in V_\lambda\) fixed by
\(\eta_s\rho_\lambda(j(\gamma_s)^{-t})\), and their
embedding is \(\phi_c=\mathcal E_{\lambda,1}(\sigma P_c)\).
Equation (CV.3) intertwines these two actions. Consequently
\[
 \begin{split}
 \phi\text{ is the image of a rational form}
 \quad\Longleftrightarrow\quad&
 \phi\text{ belongs to the complex member space, and}\\[-2pt]
 &\phi_c=\mathcal E_{\lambda,1}(Q_c)
     \text{ for some }Q_c\in\operatorname{RatSpan}(\mu)
     \text{ for every }c.
 \end{split}                                                   \tag{CV.6}
\]
To prove the reverse implication, lift \(Q_c\) by (CV.4).
The stabilizer equation and (CV.3), followed by injectivity
in CV1, force each lifted polynomial to satisfy the rational
fixed-vector equation. Thus the tuple is the required
rational form. This constructs a preimage; no averaging
over stabilizers is used.

\paragraph{CV4: extend from class representatives to all adelic points.}
Let \(G(\mathbb A_f)\) be the unitary adelic group,
\(\Gamma=G(F)\), \(U\) a supplied level subgroup, and
\(\chi:U\to\mathbb C^\times\) its character. Suppose
\(c\mapsto[y_c]\) is a bijection
\(\mathcal C\simeq\Gamma\backslash G(\mathbb A_f)/U\).
A tuple of polynomial frame functions \(\phi_c\) extends
to an adelic function with
\[
 f(\gamma y\kappa,\iota(\gamma)z)=\chi(\kappa)f(y,z)
                                                               \tag{CV.7}
\]
if and only if
\[
 y_c=\gamma y_c\kappa,\quad z=\iota(\gamma)z'
       \quad\Longrightarrow\quad
             \phi_c(z)=\chi(\kappa)\phi_c(z').                  \tag{CV.8}
\]
The necessary implication is restriction. For sufficiency,
write \(y=\gamma y_c\kappa\) and set
\[
             f(y,z)=\chi(\kappa)
                          \phi_c(\iota(\gamma^{-1})z).         \tag{CV.9}
\]
If two decompositions are used, bijectivity first gives the
same \(c\). With
\(\delta=\gamma_2^{-1}\gamma_1\) and
\(v=\kappa_1\kappa_2^{-1}\), their equality gives
\(y_c=\delta y_cv\). Equation (CV.8) compares the two
frame arguments and proves that (CV.9) is independent of
the decomposition. It also proves (CV.7).
CV1 shows that every \(f(y,\cdot)\) is still a polynomial
frame function of the required weight. This is the
substantive content of \texttt{mem\_fiC\_iff}.

\paragraph{CV5: the actual stabilizer characters agree.}
The chart reading \texttt{ReadsK} gives a matrix \(k_\beta\)
for each chart and says that \(y_ck_{\beta(c)}\) represents
the actual decorated point \(p_c\). Representation here
retains five conditions: unitarity, the carried global
lattice, the last-coordinate residues at the type primes,
the full prescribed residues at the auxiliary primes, and
the exact rational frames at the twist primes.

These five conditions are preserved by rational unitary
translation:
\[
       g\text{ represents }p
       \quad\Longrightarrow\quad
       \gamma g\text{ represents }(\gamma,1,1)p.
                                                               \tag{CV.10}
\]
Indeed
\((\gamma g)^{-1}x=g^{-1}\gamma^{-1}x\).
This identity transports the lattice and both residue
conditions. The prescribed frame matrices become
\(\gamma\) times the old matrices, which is exactly their
decorated-point action. Unitarity is preserved by
multiplication. Restricting from the common twist-prime set
\(\mathrm{SQ}\) to the particular set \(Q_a\) preserves
these conditions: \(Q_a\subset\mathrm{SQ}\) and
\(Q_a\) is disjoint from the auxiliary-prime set. The latter
follows from \(I.\mathrm{Ok}\), since the selected good
twist primes avoid the set containing the auxiliary primes.

Write a point stabilizer as \(s=(\gamma_s,\ell_s)\).
Since \(sp_c=p_c\),
\[
       (1,\ell_s)p_c=(\gamma_s^{-1},1)p_c.
\]
If \(y_c=\gamma_sy_c\kappa\), then
\[
       \gamma_s^{-1}(y_ck_{\beta(c)})
                     =y_c\kappa k_{\beta(c)}.
                                                               \tag{CV.11}
\]
By (CV.10), the right side represents
\((1,\ell_s)p_c\). Apply the character-reading clause of
\texttt{ReadsK} to this purely level action. The step
character ignores the rational unitary component, so it gives
\[
                    \chi(\kappa)=\sigma(\eta_s).              \tag{CV.12}
\]
The chart's stabilizer-range identity says precisely that
every rational \(\gamma\) with
\(\gamma\in y_cUy_c^{-1}\) comes from a point stabilizer,
with its same rational matrix. Since
\(y_c=\gamma y_c\kappa\) is equivalent to
\(\gamma=y_c\kappa^{-1}y_c^{-1}\), (CV.12) identifies
(CV.5) with (CV.8), in both directions. Applying CV4 on
each chart proves the actual complex member-space clause.

\paragraph{CV6: local bases, scalar extension and the two duals.}
For a rational prime \(p\), write \(\mathcal O_{E,(p)}\)
for the localization obtained by inverting rational integers
prime to \(p\). A matrix \(b\in\mathrm{GL}_n(E)\) is the
source's \(\operatorname{IsLocBasis}(p,H)\) exactly when
\[
                  H_{(p)}=b\mathcal O_{E,(p)}^n.              \tag{CV.13}
\]
If \(\mathfrak q\mid p\), every integer prime to \(p\) is a
unit at \(\mathfrak q\), by Bezout's identity. Localizing
(CV.13) therefore gives
\(H_{\mathfrak q}=b(\mathcal O_E)_{\mathfrak q}^n\).
Conversely, common denominators for the finitely many
coordinates justify this equality in the source's explicit
localization convention. A matrix as in (CV.13) exists for
every full lattice \(H\): over the semilocal Dedekind ring
\(\mathcal O_{E,(p)}\), a finite torsion-free module is free.
One can choose generators of its ideal summands by imposing
their valuations simultaneously at the finitely many primes
over \(p\), using CRT. This is the same rational semilocal
basis construction used in CS3.

For a coefficient prime \(w\), let
\(\mathfrak q_\tau=j_\tau^{-1}(w)\).
It lies over the same rational prime \(p_w\), since
\(j_\tau(p_w)=p_w\). Put
\[
 M_\tau=\mathcal O_L\text{-span}\,j_\tau(H),\qquad
 M_\tau^\vee=\{\ell:\ell^tx\in\mathcal O_L
                                  \text{ for every }x\in M_\tau\}.
\]
This last dual is the ordinary linear dual, with no
conjugation. If \(b\) satisfies (CV.13) for \(p_w\), then
\[
 (M_\tau)_w=j_\tau(b)R_w^n,\qquad
 (M_\tau^\vee)_w=j_\tau(b)^{-t}R_w^n.                          \tag{CV.14}
\]
For the first identity, a denominator outside
\(\mathfrak q_\tau\) maps outside \(w\), so both localized
inclusions survive scalar extension. For the second, the
pairing with the columns of \(j_\tau(b)\) identifies the
local dual. To pass from local pairings to the localized
global dual, clear the finitely many pairing denominators
on a finite generating set of \(M_\tau\).
Apply this construction to \(H=H_c^\#\).
The Hermitian dual used to obtain \(H_c^\#\) and the
ordinary dual in (CV.14) perform different operations;
both are retained.

\paragraph{CV7: localization of the transported Weyl lattice is exact.}
For any full \(H\subset E^n\), the source's
\(\mathcal D_\lambda(H)=\texttt{delAtr}(j,\lambda,H)\)
is the global polynomial lattice obtained from the ordinary
duals \(M_\tau^\vee\). At a prime \(w\), choose bases
\(g_{\tau,w}\) of these duals, and set
\[
 A_w=\rho_\lambda(g_w)
                      (\Delta_\lambda^{\mathrm{std}})_w.
\]
This does not depend on the bases: two choices differ by
matrices in \(\mathrm{GL}_n(R_w)\), which preserve the
integral Weyl lattice and its determinant twist in both
directions. To see the integrality directly, eliminate
an invertible matrix over the local ring using unit pivots.
Its elementary factors act by finite divided-power
expansions; diagonal units and permutations preserve the
integral weight decomposition. Their determinant powers
are units. This is integral stability, not a comparison
of dimensions.

By definition \(\mathcal D_\lambda(H)=\bigcap_w A_w\)
inside \(V_\lambda\), and
\[
                   \mathcal D_\lambda(H)_w=A_w.              \tag{CV.15}
\]
Here is a denominator proof of this localization step.
All the lattices \(M_\tau^\vee\) agree with the standard
ones away from finitely many primes: choose one nonzero
integer bounding each of these finitely generated full
lattices above and below the standard lattice.
Consequently \(A_w=(\Delta_\lambda^{\mathrm{std}})_w\)
away from a finite set. At each of its remaining primes,
\(A_w\) is a full \(R_w\)-lattice. Clear its finitely many
basis denominators using multipliers in \(\mathcal O_L\);
these are integral at all other primes and therefore
preserve their local modules. Their finite product gives
a nonzero common multiplier.
Thus every element of \(V_\lambda\) has one nonzero
multiple belonging to all \(A_w\). Each \(A_w\) is local
at \(w\), so LP2 proves (CV.15).

For the basis \(b\) of CV6, (CV.14)--(CV.15) give
\[
 \begin{split}
 P\in\mathcal D_\lambda(H)_w
 \quad\Longleftrightarrow\quad&
 P=\rho_\lambda\bigl((j_\tau(b)^{-t})_\tau\bigr)P_0\\
 &\text{for some }
             P_0\in(\Delta_\lambda^{\mathrm{std}})_w.
 \end{split}                                                   \tag{CV.16}
\]
Using CV1 and CV2 converts this into the precise value
criterion
\[
 \begin{split}
 P\in\mathcal D_\lambda(H)_w
 \quad\Longleftrightarrow\quad&
 \mathcal E_{\lambda,1}(\sigma P)(z)
      =\mathcal E_{\lambda,\iota(b^{-1})}(Q)(z)
                       \quad\text{for every frame }z,\\
 & Q\in\operatorname{IntSpan}_w(\mu)
                  \text{ for some }Q .
 \end{split}                                                   \tag{CV.17}
\]
The reverse direction uses polynomial injectivity from CV1.
The transpose in (CV.16) becomes exactly \(b^{-1}\) in
(CV.17), and its determinant factor is included in
\(\mathcal E\). No conjugation or determinant factor is
dropped.

\paragraph{CV8: one denominator works for all classes.}
The source's actual lattice is
\[
 \mathcal L=
 \{x\in V_{\mathrm{form}}:
                x_c\in\mathcal D_\lambda(H_c^\#)
                               \text{ for every }c\}.
                                                               \tag{CV.18}
\]
Since \(\mathcal C\) is finite,
\[
 x\in\mathcal L_w
 \quad\Longleftrightarrow\quad
 x_c\in\mathcal D_\lambda(H_c^\#)_w
                              \text{ for every }c.            \tag{CV.19}
\]
The forward implication is immediate. For the reverse,
choose \(s_c\notin w\) clearing the denominator for class
\(c\), and take \(s=\prod_c s_c\).
Then \(s\notin w\) and \(sx\in\mathcal L\).
Scalar multiplication preserves the stabilizer equations,
so this is a denominator within the rational forms space.
The source's valuation module \(\mathcal M.S(i,w)\)
is exactly \(\mathcal L_w\), because its \texttt{intW}
is the same \(R_w\). Also
\(\mathcal L=\bigcap_w\mathcal L_w\): a vector outside
\(\mathcal L\) has a proper denominator ideal contained
in a prime, contradicting membership in that localization.

\paragraph{CV9: the actual valuation-compatible chart.}
Supply the source charts \(K_\beta\), the class map
\(\beta:\mathcal C\to\mathrm{Ch}\), their \texttt{ReadsK}
identity, and \(e_Y\) satisfying \texttt{IdY}.
The source hypotheses are exactly
\[
 I.\mathrm{Ok},\quad\mathrm{Pins}(B,X),\quad
 l:I.\mathrm{Laws},\quad
 Z:\mathrm{Later}(B,I.\mathrm{cols}(l)),\quad i\le m.
                                                               \tag{CV.20}
\]
The bad-place set may be the one supplied to these
charts; no extra exclusion on \(w\) is needed.
The standard-coordinate identification of SF gives
\(\operatorname{gramN}=I\),
\(\operatorname{latCrd}(I.\mathrm{latAt}(c))=H_c\),
and \(\operatorname{dualSet}(H_c)=H_c^\#=\texttt{Lsh}_c\).
The embedding identity is the literal
\(\operatorname{embT}(D)_\tau=\sigma\circ j_\tau\),
proved from the base datum's prescribed embedding
identity, and (CV.1) identifies the weights.

CV4--CV5 give
\[
 \phi\in\mathcal M.\mathrm{Fi}(i)
 \ \Longleftrightarrow\
 \bigl(c,z\mapsto\phi_c(e_Y^{-1}z)\bigr)
       \in\operatorname{FiC}(K_\beta)
                   \quad\text{on every chart }\beta.         \tag{CV.21}
\]
CV3 gives (CV.6) with the actual
\(\mathcal M.\mathrm{emb}(i)\), \(e_Y\), and
\(\operatorname{RatSpan}(D,\mathcal M.\mathrm{lam}(i))\).
Finally CV7--CV8 give, for every coefficient prime \(w\),
\[
 \begin{split}
 x\in\mathcal M.S(i,w)
 \quad\Longleftrightarrow\quad
 &\text{for every }c\text{ and every }
    b\in\operatorname{IsLocBasis}(p_w,H_c^\#),\\
 &\text{there is }Q\in\operatorname{IntSpan}_w(\mu)
    \text{ such that, for every }z,\\[-2pt]
 &\mathcal M.\operatorname{emb}(i)(x)_c(z)
       =\mathcal E_{\lambda,\iota(b^{-1})}(Q)(e_Yz).
 \end{split}                                                   \tag{CV.22}
\]
The reverse implication is not vacuous: CV6 supplies at
least one such \(b\) for each class and residue prime.
These are precisely the three clauses of
\texttt{ValuationCompatibleChart}, with
\(\texttt{okW}\equiv\mathrm{True}\).
They prove \texttt{TC.val\_tied} and its exported
\texttt{hermitianChart\_complex\_rational\_integral\_forms}
statement. The all-prime integral criterion is now proved
for the supplied charts, rather than placed among their
assumed properties.

\paragraph{CV10: the smaller terminal row space uses the same proof.}
For \(\texttt{clsP\_tie}\), take \(n'=N_m-1\), the actual
row step \(I.\mathrm{stX}(\mu_0)\), and its supplied row
class chart with \texttt{ReadsP}. Its untwisted coefficient
weight is \(\mu'=\texttt{sgW}\), and its frame weight is
\[
                  \lambda'=\mu'-\frac{m(t+1)}2(1,\ldots,1).
                                                               \tag{CV.23}
\]
The equality of the two source definitions follows from
the same oddness calculation as (CV.1). Thus both
\(V_{\lambda'}=V_{\mu'}\) and
\(\Delta_{\lambda'}^{\mathrm{std}}
 =\Delta_{\mu'}^{\mathrm{std}}\), while the frame
determinant exponent stays that of \(\lambda'\).
At this row step the type-prime decoration set is empty.
The decorated lattice \texttt{Lsh} therefore equals its
carried lattice \(N_c\) exactly.

CV4--CV5 prove the terminal \texttt{ptC} clause using the
row chart's same stabilizer matrices and character.
CV3 proves its \texttt{rat} clause.
CV6--CV8, followed by intersection over all coefficient
primes, prove its \texttt{lat} clause with \(N_c\) in place
of \(H_c^\#\). The frame identification is the identity
on matrix entries, as in \texttt{clsP\_tie}.
The supplied terminal chart and its remaining construction are tracked in FC; this
corollary supplies its value clauses. Its
\texttt{ipP} and \texttt{NormPTie} clauses require the
separate normalized Fock and Hermitian metric argument
and are not consequences of (CV.21)--(CV.22).

\paragraph{Boundary of this construction.}
The proof constructs rational preimages and proves the
exact integral coefficient criterion for the actual
step lattices. It retains the supplied true column datum,
its \texttt{Ok/Pins/Laws/Later} witnesses, and the chart
reading identities used in CV5. FC describes their
simultaneous chart and carrier assembly. CV does not
infer those inputs from their names, identify an
integral Weyl lattice with its dual-Weyl enlargement,
or infer a metric identity from the lattice equality.
